% Academic revision prepared from songlines_unified(2).tex.
% Revision goals: conventional problem-method-results-scope structure;
% cautious claim language; terminology consistency; preservation of all
% equations, tables, labels, references, and reported numerical results.
\documentclass{article}

\usepackage{neurips_2026}

\usepackage[utf8]{inputenc}
\usepackage[T1]{fontenc}
\usepackage{hyperref}
\hypersetup{hidelinks}
\usepackage{url}
\usepackage{booktabs}
\usepackage{amsfonts}
\usepackage{amssymb}
\usepackage{amsmath}
\usepackage{nicefrac}
\usepackage{microtype}
\usepackage{xcolor}
\usepackage{graphicx}
\graphicspath{{./}}
\usepackage{array}
\usepackage{enumitem}
\usepackage{placeins}
\usepackage{tikz}
\usetikzlibrary{arrows.meta,positioning}

\newcommand{\method}[1]{\texttt{#1}}
\newcommand{\wstar}{W^{\star}}
\newcommand{\rwstar}{RW^{\star}}
\newcommand{\mstar}{M^{\star}}
\newcommand{\cstar}{C^{\star}}
\newcommand{\phiroute}{\varphi_{\mathrm{route}}}

\setcounter{tocdepth}{2}

\title{The Song, Not the Pin: Provenance-Conditioned Route Transfer across Private Frames}

\begin{document}
\maketitle
\begin{abstract}
Collective-memory systems allow an agent to act on evidence that was
acquired by another agent, yet conventional task and stage metrics do
not identify when a commitment is supported primarily by foreign
memory. We introduce a provenance-conditioned extension of Q/R/M/C.
For each materialized target, the method records the foreign share
$\varphi$ of the evidence mass used by the memory merge and freezes
this value at commitment time. A semantic-warp event is then defined
as a materialization whose supporting evidence exceeds a specified
foreignness threshold. This annotation permits three analyses: direct
completion stratified by provenance, deterministic counterfactual
replay with foreign evidence removed, and a closed-form admissibility
horizon derived from the trust--staleness gate. On a 2,400-episode
scarcity benchmark, foreign-evidence locks complete substantially less
often than own-evidence locks ($0.004$--$0.020$ versus
$0.22$--$0.41$), but this direct comparison understates their value.
Lock-chain analysis shows that foreign evidence frequently transports
an agent into self-confirmation range: up to $51\%$ of subsequent
own-evidence completions were preceded by a foreign-evidence transport
chain (a descriptive chain statistic; the causal contribution is
estimated separately by masked replay). The closed-form admissibility
boundary derived from the trust--staleness gate predicts the latest
admissible intention time exactly on the registered $6/6$ hold-out and
on a second, parameter-disjoint $75$-cell hold-out ($75/75$ exact,
including delayed snapshot delivery, whose breakpoint shift the closed
form predicts tick-for-tick) --- a contract-level prediction validated
against the implementation, not an empirical law of the environment. The same analysis identifies target contention, rather
than incorrect evidence, as the principal completion cost of fast
sharing --- a diagnosis confirmed by intervention: restoring
coordination alone (unique target assignment, evidence untouched)
restores foreign-evidence completion from $0.000$ to $0.48$--$0.56$,
whereas perfect evidence without coordination does not help. We
therefore introduce a decentralized reservation and rollback protocol
that increases success at every tested fast cadence, exceeds the
strongest fixed-cadence baseline ($0.614$--$0.620$ versus $0.583$),
and, against five standard decentralized coordination baselines on
equal information, is best or statistically tied on success while
uniquely minimizing duplicate commitments at every cadence. A controlled LLM experiment further demonstrates
cross-session use of foreign memory. Extending provenance from target
records to route edges then reveals a structural boundary between
place-only and route-level transfer: place-only guidance fails at the
first obstructing structure and can perform worse than unguided
exploration, whereas transferred routes remain executable. An
edge-specific rupture boundary, derived from the same gate, predicts
all $110$ tested validity regimes and all $21$ intermediate failure
indices exactly, and a traveler on a foreign route preserves the
witness's measured risk profile up to the predicted rupture. Semantic landmark constellations recover
correspondence across translations and quarter-turn rotations in $80/80$
tests, remain fail-closed across a $5{,}280$-trial stress surface of
tag noise, landmark deletion, and constructed ambiguity, and expose
one systematic fail-open channel --- silent vocabulary mismatch
($2.5$--$5\%$ wrong-transform acceptance under permuted tag alphabets)
--- which makes the shared-ontology assumption explicit. Landmarks and edges provide
complementary identity and transport functions; their composition
reaches an unseen target at $6.4\%$ of the full-snapshot bit cost before
protocol metadata. These results indicate that the
value of shared memory depends on the provenance, temporal validity,
and coordination authority attached to a record, not only on its
semantic content. The scope is limited to controlled grid and
continuous simulations and a small text-world demonstration.
\par\medskip\noindent\textbf{We claim.} In the evaluated class of
partially observed navigation settings with structural obstacles ---
where the recipient cannot independently reconstruct a feasible path
from a destination record --- foreign navigation memory becomes
executable only when referential identity, relational transport,
evidence provenance, temporal validity, and coordination authority are
preserved together. In this class, a target without a route is
operationally incomplete; a route without a valid cross-frame
correspondence is unsafe; and foreign evidence may guide bounded,
reversible transport but should not acquire durable action authority
without a validity and coordination contract. The executable unit of
transferable navigation memory here is therefore not a coordinate
target but a provenance-conditioned relational route whose landmarks
establish identity, whose edges provide transport, and whose validity
contract determines how far and under what conditions another agent
may execute it.
\end{abstract}


\FloatBarrier
\section{Introduction}
\label{p2:sec:intro}

A shared memory architecture is usually described in terms of the
information available to each agent. This description is incomplete.
Two agents may retrieve the same target record, but the epistemic
basis of their commitments can be different: one may rely on direct
observation, whereas the other may rely entirely on a peer's report.
End-to-end success, retrieval accuracy, and the original Q/R/M/C
rates do not distinguish these cases because they do not condition
materialization and completion on evidence provenance.

This paper introduces a provenance-conditioned measurement layer for
collective semantic memory. Every contribution to a merged record
retains its source, support, confidence, trust, and age. For a selected
candidate, the method computes the foreign share $\varphi$ using the
same weights that produced the merge. The value is frozen when the
agent commits to the target so that later self-confirmation does not
retroactively change the provenance of the decision. A semantic-warp
event is therefore not an additional stage, but a stratification of
materialization and completion according to the source of the
supporting evidence.

The analysis separates three questions. First, how often do
foreign-evidence commitments complete directly? Second, what is the
counterfactual effect of foreign memory when its contribution is
removed under an otherwise identical replay? Third, for how long may
foreign evidence remain action-authoritative under explicit trust and
staleness rules? The symbolic memory representation makes all three
questions operational: the merge is source-decomposable, the simulator
supports paired deterministic replays, and the exponential staleness
gate yields a closed-form horizon.

The results show that direct completion is only one component of
foreign-memory value. Under target scarcity, foreign locks rarely
complete because several agents can commit to the same resource.
Nevertheless, the same lock often guides the recipient to a region in
which it can obtain its own confirming evidence. We therefore treat
foreign memory as a transport mechanism as well as a destination
proposal. This distinction explains why an architecture can receive a
positive counterfactual benefit from foreign evidence even when
$P(C^\star\mid W^\star)$ is small.

The observed failure is primarily coordinative. Fast exchange
increases the number of commitments based on the same peer evidence,
which improves materialization but increases duplicate target claims.
To address this mechanism, we introduce a minimal decentralized
protocol consisting of optimistic commitment, reservation messages,
and rollback with backoff. The protocol does not alter the semantic
content of the memory; it regulates when a shared record is permitted
to determine action.

The contributions of this paper are as follows:
\begin{enumerate}[leftmargin=20pt,itemsep=2pt,topsep=2pt]
\item We define provenance-conditioned target commitments and a
deterministic counterfactual measure of foreign-evidence contribution.
\item We derive and validate a trust--staleness admissibility horizon.
\item We identify downstream transport and target contention as
distinct consequences of foreign evidence and introduce reservation and
rollback.
\item We extend provenance from target records to route edges and
define strict route warp.
\item We derive an edge-specific rupture boundary and validate the
predicted failure edge.
\item We introduce meaning-based correspondence across private frames
and demonstrate fail-closed behaviour under ambiguity.
\item We separate identity-bearing landmarks from transport-bearing
edges and quantify their combined storage cost.
\end{enumerate}

The claims are limited to the tested memory rules and environments.
The admissibility boundary is exact under the specified exponential gate; it is
not asserted for arbitrary learned memory policies. Similarly, the LLM
study demonstrates measurability and cross-session transfer, but does
not constitute a broad benchmark of hallucination-resistant agent
collectives.

\paragraph{Structure of the paper.} The paper has two movements on a
single object. The first (Sections~\ref{p2:sec:setting}--\ref{p2:sec:llm})
treats foreign evidence about an isolated \emph{target}: it defines the
provenance-conditioned commitment, measures its counterfactual value,
derives the trust--staleness admissibility horizon, and adds the
reservation-and-rollback coordination contract. The second movement,
beginning at Section~\ref{p2:sec:bridge} and developed in
Sections~\ref{p3:sec:routes-lead}--\ref{p3:sec:identity}, shows that the transferable unit is
not a target but a \emph{relational route}: landmarks establish
cross-frame identity, edges provide transport, edge-level provenance
sets a rupture boundary, and only the composition reaches an unseen
target. The two movements share one contract --- identity, transport,
provenance, validity, and coordination authority preserved together.


\section{Setting}
\label{p2:sec:setting}

\paragraph{Substrate and task.} $N$ agents explore a $12\times10$ grid
with $T$ semantic water targets ($N>T$: scarcity) and hazards. Each
agent maintains symbolic place memory (tags, confidences, support,
freshness). Four sharing architectures over the identical substrate
and planner: \method{independent} (no communication),
\method{shared} (single event bus), \method{centralized} (consensus
aggregator), \method{peer} (broadcast every $K$ ticks, asymmetric
trust), plus \method{CSM} (peer $+$ explicit merge/trust/staleness
rules with decay rate $\alpha{=}0.05$ and inclusion threshold
$\tau{=}0.30$).

\paragraph{Q/R/M/C instrumentation.} Per tick and per agent we log
query formation (Q), retrieval satisfaction (R), target
materialization as the planner's lock ($\mstar$), and completion
($\cstar$). Episode-level conditional rates follow the Q/R/M/C event contract defined in the companion Q/R/M/C evaluation study \citep{songlines1}. All warp measurements below are annotations of these
$\mstar$ events; no planner or memory invariant is modified.


\section{The warp event}
\label{p2:sec:warp-def}

\paragraph{Provenance: the foreign evidence mass ratio.} Every
evidence unit entering a merge carries its source. For the candidate
$m$ the planner locks,
\[
\varphi(m) \;=\;
\frac{\sum_{e\in\mathrm{supp}(m),\,\pi(e)\neq\mathrm{self}} w(e)}
     {\sum_{e\in\mathrm{supp}(m)} w(e)},
\]
where $w(e)$ are \emph{the merge's own weights}: $\text{trust}\cdot
\log(1{+}\text{support})$ for peer/centralized consensus,
$\text{trust}\cdot e^{-\alpha\,\text{age}}\cdot\text{conf}$ for CSM,
observation masses for the shared bus. For peer and centralized
architectures $\varphi$ is reconstructed exactly from the
per-contribution fields the merge already exports; the measurement
layer reads state and never writes it.

\paragraph{What $\varphi$ is and is not.} $\varphi$ is a
\emph{descriptive} quantity --- the share of merge weight contributed
by other agents --- not a causal contribution. A commitment with
$\varphi{=}0.8$ is not thereby ``80\% caused'' by foreign evidence: a
small own share can be pivotal (without it the candidate would not
clear the threshold), and a large foreign share can be redundant. We
therefore keep three quantities apart. \emph{Mass}: $\varphi(m)$, read
off the merge. \emph{Pivotality}: $\Pi_i = \mathbf{1}[a_i(M^{\mathrm{self+foreign}})
\neq a_i(M^{\mathrm{self}})]$ --- whether removing foreign evidence
changes the agent's commitment. \emph{Outcome effect}:
$\Delta_i = Y_i(M^{\mathrm{self+foreign}}) - Y_i(M^{\mathrm{self}})$
--- the change in the episode outcome. $\varphi$ stratifies the logs
(this section); $\Pi$ and $\Delta$ are estimated by the masked
deterministic replay of \S\ref{p2:sec:gain}.

\paragraph{Freezing rule.} $\varphi$ is fixed at the moment of the
lock. If the agent later confirms the place with its own eyes,
$\varphi$ is \emph{not} recomputed --- otherwise every successful warp
would dilute itself on approach and $P(\cstar|\wstar)$ would be
unmeasurable by construction.

\paragraph{Warp events and strata.} $\wstar_\theta := \mstar \wedge
(\varphi\ge\theta)$ with $\theta{=}1$ (strict: never self-observed,
cross-checked against an observation ledger) and $\theta{=}0.8$
(soft). Warp share $P(\wstar|\mstar)$ is a structural descriptor
(identically $0$ for \method{independent}); warp completion
$P(\cstar|\wstar)$ is the new quantity. Under the Q/R/M/C event
contract defined in the companion Q/R/M/C evaluation study \citep{songlines1}, warp is an exact
stratification of the materialization event $\mstar$ rather than an
additional stage: identifiability of Q/R/M/C is untouched.

\paragraph{Sensitivity to the soft threshold.} The soft threshold
$\theta{=}0.8$ is a reporting convention, not a load-bearing constant.
Sweeping $\theta \in \{0.5, 0.6, 0.7, 0.8, 0.9, 0.95\}$ on the same
lock-level logs (artifact \texttt{tmp/warp/phi\_sensitivity.json})
leaves the headline stratification intact at every threshold and both
tested cadences: at $K{=}2$, $P(\cstar|\wstar_\theta)$ stays in
$[0.004, 0.005]$ while the own stratum stays in $[0.26, 0.50]$; at
$K{=}8$, $[0.014, 0.027]$ against $[0.41, 0.48]$. The foreign stratum
completes an order of magnitude less often than the own stratum at
\emph{every} $\theta$; only the warp share moves (from $0.65$ at
$\theta{=}0.5$, $K{=}2$ down to $0.09$ at $\theta{\ge}0.9$, $K{=}8$),
as it must when the stratum boundary shifts.

\paragraph{Regression anchor.} The known misled-agent episode of the
architecture comparison (agent A locks $(3,8)$ known only from C; C
claims it first; A strands) is annotated by the instrumentation as
$\varphi=1.0$ strict $\wstar$ with collision --- the pathology the
framework must see, and the protocol of \S\ref{p2:sec:drive} must fix.


\section{Direct and downstream effects of foreign evidence}
\label{p2:sec:raw}

\paragraph{Design.} $2{,}400$ episodes: $(N,T)\in\{(3,2),(5,3),(8,5),
(4,4),(6,2)\}$ $\times$ three layouts $\times$
$\{\text{peer-}K2,\text{peer-}K8,\text{CSM-}K8,\text{shared}\}$
$\times$ 20 seeds $\times$ \{full, foreign-masked\} with identical
seeds. Unit of analysis: lock event; CIs by episode-cluster bootstrap.

\begin{figure}[t]
\centering
\includegraphics[width=0.72\linewidth]{figures/fig_warp_strata}
\caption{$\varphi$-stratified completion vs.\ broadcast cadence $K$
(dedicated peer sweep: scarcity cells, random layouts, 20 seeds,
episode-cluster bootstrap CIs). As $K$ shrinks, the warp share
$P(\wstar|\mstar)$ grows from $0.11$ to $0.55$ while
$P(\cstar|\wstar)$ stays below $0.05$ and the own-stratum rate stays
flat: the completion collapse of fast sharing is concentrated in the
$\wstar$ stratum, giving the Q/R/M/C bottleneck shift reported in the companion Q/R/M/C evaluation study \citep{songlines1} its
mechanistic reading. The share uptick at $K{=}16$ (broadcasts are
rare, but agents also discover little on their own) is reported as
observed.}
\label{p2:fig:strata}
\end{figure}

Figure~\ref{p2:fig:strata} shows the cadence view on a dedicated peer
sweep ($K\in\{1,2,4,8,16\}$): warp share is what fast sharing buys,
and it is precisely the stratum that does not complete.
Table~\ref{p2:tab:hw1} gives the architecture view.

\begin{table}[h]
\centering
\caption{Provenance strata of completion under scarcity ($N>T$), full
runs. Warp locks complete an order of magnitude less often than
own-evidence locks; all CI pairs are disjoint. (H-W1 confirmed.)}
\label{p2:tab:hw1}
\small
\begin{tabular}{lcccc}
\toprule
Architecture & $P(\cstar|\wstar_{0.8})$ & $n_{\wstar}$ & $P(\cstar|\mstar,\text{own})$ & $n_{\text{own}}$ \\
\midrule
peer-$K{=}2$ & 0.004 [0.001, 0.008] & 1{,}437 & 0.223 [0.186, 0.275] & 2{,}738 \\
peer-$K{=}8$ & 0.020 [0.000, 0.038] & 149    & 0.389 [0.329, 0.446] & 1{,}734 \\
CSM-$K{=}8$  & 0.017 [0.006, 0.044] & 637    & 0.414 [0.368, 0.453] & 1{,}684 \\
shared       & 0.004 [0.002, 0.008] & 2{,}099 & 0.279 [0.237, 0.332] & 2{,}250 \\
\bottomrule
\end{tabular}
\end{table}

A caveat on reading Table~\ref{p2:tab:hw1}: the strata are
observational, and own- and foreign-evidence locks do not arise under
matched conditions. On the pooled lock-level logs ($23{,}237$ events;
artifact \texttt{tmp/warp/phi\_matched\_analysis.txt}) the covariate
imbalance is severe: own locks are made $1.15$ cells from the target
on average, foreign locks $6.26$; co-lock pressure is $0.21$ vs
$1.76$. A covariate-matched comparison is possible only in the small
overlap region (foreign locks made within $2$ cells): there the
own--foreign completion gap narrows from $0.513$ unadjusted to
$0.38$--$0.45$ within (distance $\times$ co-lock) strata --- covariates
explain roughly a quarter of the raw gap, and the remainder largely
reflects the freezing rule's bookkeeping, whereby an approached
foreign lock is typically dropped and re-locked as own evidence, so
its completion is booked in the own stratum
(\S\ref{p2:sec:taxi}). The near-disjointness of the two covariate
distributions is itself the transport story: foreign locks are the
locks made from far away. The causal quantity is estimated by the
paired masked replay of \S\ref{p2:sec:gain}. The same logs give the
continuous completion curve over $\varphi$: mass is nearly bimodal
(locks sit at $\varphi{=}0$ or $\varphi \gtrsim 0.6$), and completion
falls off between the $[0.4,0.6)$ bin ($0.07$--$0.22$ per
architecture) and the $[0.6,0.8)$ bin ($\le 0.012$), so the reported
strata do not depend on the exact placement of the soft threshold.

\subsection{Downstream transport through lock chains}
\label{p2:sec:taxi}

Direct completion understates the value of foreign evidence. A
\emph{warp-assisted own completion} is a completed own-evidence lock
preceded (same agent, same episode) by a soft-$\wstar$ lock that was
dropped: the warp failed as a destination but succeeded as a ride.

\begin{table}[h]
\centering
\caption{Lock-chain analysis (random layouts, full runs). The first
$\wstar$ lock in a chain is made from ${\approx}8$--$10$ cells away;
the completing own lock from $2.0$ cells --- exactly the observation
radius, i.e.\ the self-confirmation point.}
\label{p2:tab:taxi}
\small
\begin{tabular}{lccc}
\toprule
Arm & own $\cstar$ & warp-assisted (share) & $r(\text{first }\wstar) \to r(\text{own})$ \\
\midrule
shared          & 248 & 127 (\textbf{51.2\%}) & $9.5 \to 2.0$ \\
peer-$K{=}2$    & 223 & 99 (\textbf{44.4\%})  & $9.7 \to 2.0$ \\
CSM-$K{=}8$     & 279 & 13 (4.7\%)            & $5.8 \to 2.0$ \\
peer-$K{=}8$    & 246 & 7 (2.8\%)             & $8.4 \to 2.0$ \\
\bottomrule
\end{tabular}
\end{table}

This resolves an apparent paradox in the counterfactual data: the
shared bus at moderate scarcity gains $+10.8$ ticks $[+2.2,+20.9]$
from foreign evidence although its direct $P(\cstar|\wstar)$ is
$0.004$. The gain flows through chains; the freezing rule books the
completion in the own-evidence stratum, and the chain analysis
surfaces the transport chain as a separate measured quantity. The
chain share is a \emph{descriptive precedence} statistic --- up to
$51\%$ of own-evidence completions \emph{were preceded by} a
foreign-evidence transport chain --- not a causal attribution: the
sequence does not by itself show the agent would have missed
self-confirmation range without the foreign lock. The causal
counterpart is the masked-replay gain of \S\ref{p2:sec:gain}, which
removes the foreign evidence and measures the outcome change under
the same seed.

\subsection{Counterfactual warp gain}
\label{p2:sec:gain}

Two estimands must be distinguished, because in a multi-agent system
foreign evidence is an interference channel and no SUTVA-style
no-interference condition holds. The \emph{joint system effect}
masks foreign evidence at merge for \emph{all} agents and replays the
episode under the same seed and random stream: it answers ``what does
sharing as a system-level mechanism contribute?''. The
\emph{individual conditional effect} masks foreign evidence only in
agent $i$'s memory view while the other agents share normally: it
answers ``what does received foreign evidence contribute to this
agent, spillovers included?''. Both are paired deterministic system
replays under the registered simulator and intervention policy; after
the first divergent action the two worlds differ, so each estimand is
a total-system counterfactual, not a per-decision mediation effect.
We report both. Masking foreign evidence at merge under the same seed
gives $\mathrm{WG} = t_{\text{cens}}(\text{masked}) -
t_{\text{cens}}(\text{full})$ per episode. The sign map over (layout,
scarcity $\rho{=}N/M$) is regime-dependent: positive for the shared
bus at unpredictable layouts and moderate scarcity ($+10.8$ at
$\rho{=}1.5$), negative for fast peer sharing (peer-$K2$ random:
$-7.1$ $[-13.5,-1.3]$), ${\approx}0$ at hard scarcity. A per-agent
mask (one agent masked, others share normally) confirms the sign at
the individual level: $\mathrm{WG}_i = -10.25$ $[-20.91,-0.16]$ ticks
on the fast-share scarcity cell, with spillover on unmasked agents
${\approx}0$ ($-0.62$ $[-3.06,+1.91]$). Raw warp \emph{hurts} its
recipient precisely where sharing is fastest --- the collision cost
exceeds the information value.

\paragraph{Exploratory: collision pressure.} Co-lock pressure (other
agents locked on the same target) correlates with warp failure with
the predicted sign but negligible magnitude (Spearman $+0.033$,
$p{=}0.018$, $n{=}5{,}117$); we report it as exploratory rather than
as a confirmed predictor.

\paragraph{Contention, established by intervention.} The claim that
target contention --- not incorrect evidence --- is the principal
completion cost of fast sharing is closed by a $2{\times}2$
evidence-quality $\times$ coordination intervention grid plus an
abundance control ($800$ episodes, $20$ seeds, identical seeds across
modes; artifact \texttt{tmp/cluster/contention\_interv/}). With
evidence untouched and coordination fixed by a unique-assignment
oracle, foreign-evidence completion is \emph{restored}:
$P(\cstar|\wstar_{0.8})$ rises from $0.000$ to $0.478$ at $K{=}1$ and
$0.559$ at $K{=}4$, with duplicate locking eliminated. Removing
scarcity ($T{=}N$) lifts success from $0.60$ to $0.80$--$0.83$
without any coordination. Conversely, perfect evidence \emph{without}
coordination (oracle retrieval, no protocol) leaves success within
noise of baseline and duplicate locking high. The completion penalty
of foreign evidence is therefore a coordination failure, not an
information failure --- the direct interventional counterpart of the
descriptive co-lock statistics above.


\section{The warp admissibility boundary}
\label{p2:sec:law}

CSM includes foreign evidence with weight $\tau_i\, e^{-\alpha\,
\text{age}}\, c$ (trust $\tau_i$, confidence $c{=}0.95$) and inclusion
threshold $\tau{=}0.30$. Foreign evidence therefore stops forming
intention at
\[
\mathrm{age}_{\max}(\tau_i) \;=\; \tfrac{1}{\alpha}\,
\ln\!\big(\tau_i\, c / \tau\big), \qquad \text{defined for } \tau_i c > \tau .
\]

\paragraph{Witness--traveller protocol.} A witness observes water once
and goes silent; a traveller who has never seen it receives the
snapshot at a parameterised age. Every lock the traveller makes is a
strict $\wstar$ by construction. Figure~\ref{p2:fig:law} summarises both
experiments.

\begin{figure}[t]
\centering
\includegraphics[width=\linewidth]{figures/fig_warp_law}
\caption{The warp admissibility boundary. \textbf{(a)} Foreign-evidence merge
weight under trust$\times$staleness decay; $\times$ marks the
empirically measured last age at which the traveller's query still
exposes the target --- each sits on its curve's crossing of the
inclusion threshold $\tau$. Fixed-cadence peer memory has no staleness
term (dashed): its gate never closes; the architectures differ in the
\emph{shape} of the curve. \textbf{(b)} Predicted vs.\ empirical
completion breakpoints in full navigation episodes: original grid
(circles; age sampled in steps of 2, hence points at most one grid
step below the diagonal) and the registered hold-out (stars; age step
1, predictions written to disk before execution, $6/6$ exact,
including the never-cell at $(-1,-1)$ and the cell where the pruning
horizon undercuts the exponential gate).}
\label{p2:fig:law}
\end{figure}

\begin{table}[h]
\centering
\caption{Gate curve: largest evidence age at which the traveller's
query still exposes the target. Fixed-cadence peer memory has no
staleness term: the gate never closes (checked to age 120 at all
trusts) --- the architectures differ in the \emph{shape} of the curve,
not just its level.}
\label{p2:tab:gate}
\small
\begin{tabular}{lcc}
\toprule
trust $\tau_i$ & empirical $\mathrm{age}_{\max}$ & predicted \\
\midrule
1.00 & 23 & 23.05 \\
0.80 & 18 & 18.59 \\
0.60 & 12 & 12.84 \\
0.40 & 4  & 4.73 \\
0.25 & gate closed & 0 (trust-flip) \\
\bottomrule
\end{tabular}
\end{table}

\paragraph{Completion breakpoint.} In full navigation episodes the
feasibility inequality acquires one discrete refinement: the traveller
self-confirms the target once it is within the observation radius, so
the foreign gate only has to survive until tick $t_{\text{gate}} = d -
r_{\text{obs}} - 1$, giving breakpoint $\mathrm{bp} =
\mathrm{age}_{\max} - t_{\text{gate}}$; snapshot pruning executes on
broadcast ticks and adds the bound $10K - t_{\text{bcast}}$. All
constants of the boundary are unchanged. All 11 CSM cells of the original
grid match within sampling resolution; completion curves are monotone
steps; the trust-flip zeroes $\wstar$ events discretely.

\begin{table}[h]
\centering
\caption{Registered-prediction hold-out: six cells absent from the
original grid, predictions written to disk \emph{before} any episode,
age grid step 1. All six are exact integer matches, including the
never-cell and the cell where the pruning horizon must undercut the
exponential gate.}
\label{p2:tab:holdout}
\small
\begin{tabular}{lccc}
\toprule
Cell & binding bound & predicted & empirical \\
\midrule
$K{=}8$, $\tau_i{=}0.9$, $d{=}9$  & gate    & 14 & \textbf{14} \\
$K{=}8$, $\tau_i{=}0.7$, $d{=}9$  & gate    & 9  & \textbf{9}  \\
$K{=}8$, $\tau_i{=}0.5$, $d{=}15$ & gate    & never & \textbf{never} \\
$K{=}4$, $\tau_i{=}1.0$, $d{=}6$  & gate    & 20 & \textbf{20} \\
$K{=}2$, $\tau_i{=}1.0$, $d{=}12$ & \emph{pruning} & 12 & \textbf{12} \\
$K{=}2$, $\tau_i{=}0.8$, $d{=}8$  & gate    & 13 & \textbf{13} \\
\bottomrule
\end{tabular}
\end{table}

The boundary therefore predicts the held-out horizons exactly at tick resolution: the CSM rules
really are the mechanism that forms (and un-forms) intention. It also
reinterprets the companion result \citep{songlines3} that CSM dominates fixed-cadence
peers: trust$\times$staleness is a \emph{warp gate} --- CSM does not
share better, it \emph{refuses to warp on a stale or untrusted song}.

\paragraph{Parameter-disjoint hold-out ($75$ cells).} Because six
cells are a thin validation, the boundary was re-tested on a second
registered hold-out in which \emph{every} gate parameter takes values
never used in any prior experiment: trust
$\tau_i \in \{0.35, \dots, 0.95\}$, decay
$\alpha \in \{0.03, 0.15, 0.25\}$ (all prior runs used $0.05$),
inclusion threshold $\tau \in \{0.2, 0.4\}$, confidence
$c \in \{0.9, 0.99\}$, distances $d \in \{5,\dots,17\}$, cadences
$K \in \{2,4,8\}$, plus a new mechanism axis --- snapshot delivery
delayed by $L \in \{0,2,5\}$ ticks, for which the closed form
predicts a breakpoint shift of exactly $L$. Predictions for all $75$
cells (including $12$ never-cells and $4$ cells where the pruning
horizon must beat the exponential gate) were written to disk with a
UTC timestamp before any episode ran. Result: $\mathbf{75/75}$ exact
integer matches, $0$ off-by-one, $0$ misses; the delay shift is
$L$ tick-for-tick; all completion curves are monotone steps (artifact
\texttt{tmp/cluster/warp/boundary\_holdout\_v2/}). The boundary is a
contract-level prediction, and the contract now holds across the
parameter space, not at six points.

\paragraph{Portability to continuous dynamics.} All three registered
predictions survive on a continuous VMAS substrate through the same
continuous$\to$grid bridge used by the portability test of the companion Q/R/M/C study \citep{songlines1}:
(i) the provenance strata stay an order of magnitude apart
($P(\cstar|\wstar)=0.012$ $[0.000,0.042]$ vs.\ $0.562$
$[0.457,0.709]$ at $K{=}1$); (ii) warp share falls monotonically with
cadence ($0.723\to0.366\to0.000$ over $K{=}1/4/16$); (iii) the boundary
becomes \emph{metric}: with the travel-to-sight time calibrated once
on a fresh-evidence episode, the predicted breakpoints match within
the sampling grid, and all cells beyond the metric warp radius
($v\cdot\mathrm{age}_{\max}$ shorter than the distance) are predicted
\emph{never} and observed \emph{never}. One genuinely continuous
finding: inertia is a physical warp-leak channel --- an agent can
coast into sighting range after the gate closes --- removed by a
no-lock-then-brake controller that restores the grid semantics.


\section{The Warp Drive protocol}
\label{p2:sec:drive}

Warp collision --- locking on foreign evidence for a target someone
else will claim --- is the dominant failure of fast sharing. Warp
Drive adds three primitives, all decentralized:

\begin{enumerate}[leftmargin=20pt,itemsep=2pt,topsep=2pt]
\item \textbf{Optimistic warp lock} --- locks on foreign evidence
  remain allowed immediately (planner unchanged).
\item \textbf{Warp reservation} --- at the owner's next broadcast tick
  the reservation rides on the ordinary snapshot (same cadence $K$, no
  coordinator, no extra channel); receivers suppress reserved targets
  in their queries. Priority: (lock tick, distance at lock, agent id).
\item \textbf{Anti-$\mstar$ rollback} (Time Warp
  semantics~\cite{jefferson1985}) --- an agent holding a lock
  dominated by a foreign reservation retracts it: the planner
  re-queries, the target enters exponential backoff
  ($4\to8\to\cdots\to32$ ticks) with an agent-id tie-break against
  oscillation. The event log stays append-only.
\end{enumerate}

\paragraph{Decentralisation check.} Driving the protocol by direct
per-agent calls --- no runtime, no environment --- produces the same
retraction decisions from each agent's own inbox alone.

\paragraph{The authority ladder.} The protocol makes explicit a
distinction that runs through this paper and its companion
\citep{songlines3}: a foreign record passes through four authority
states. It is \emph{observed} (someone's local evidence);
\emph{received} (present in the recipient's merged view but inert);
\emph{provisionally executable} --- permitted to guide bounded,
reversible transport under an explicit validity contract
(staleness-gated, reservation-checked, rollback-armed, cut off at the
rupture boundary); and \emph{admitted}, with durable authority over
future decisions, which in the companion study requires validation on
the receiver's own visit. Everything in this paper operates at the
third state: a warp lock is a provisional, revocable commitment
(optimistic lock plus rollback), never a durable admission. The
ladder also dissolves an apparent circularity between the two papers
--- receiver-side validation requires reaching the place, while
reaching the place seems to require acting on the unvalidated record.
Provisional transport authority is precisely the licensed middle: the
foreign record may ferry the agent into self-confirmation range
(\S\ref{p2:sec:taxi}) under bounded risk, and admission --- the
durable upgrade --- happens only on arrival.

\paragraph{Misled-A anchor.} With Warp Drive the pathological lock is
resolved \emph{preventively}: C's reservation arrives before A ever
locks $(3,8)$; A completes on another target. In the sweep, rollback
latency of dominated locks is finite and ${\approx}K$: $1.0/1.5/2.7$
ticks at $K{=}1/2/4$ (baseline: until end of episode).

\paragraph{Against standard coordination baselines.} Reservation is a
natural remedy for duplicate claims, so Warp Drive is also compared
against five decentralized alternatives operating on the same
broadcast information at the same cadences ($2{,}520$ episodes, $20$
seeds, identical seeds across arms; per-arm message bytes accounted,
including the reservation metadata; artifact
\texttt{tmp/cluster/coord\_baselines/}): random tie-break (transitive
per-target priority), nearest-agent-wins, decentralized greedy
assignment, exponential backoff without reservations, and a soft
occupancy penalty. The honest summary: on \emph{success rate} Warp
Drive is best or statistically tied at every cadence ($0.611$--$0.620$
vs.\ no-protocol $0.570$--$0.587$), and single alternatives reach it
at single cadences (backoff $0.627$ at $K{=}1$; greedy
$0.612$--$0.615$) --- no arm dominates it, but reservation is not
uniquely capable of recovering the success loss. What distinguishes
Warp Drive is the combination: the lowest duplicate-commitment count
at every cadence ($2.8$--$5.2$ events; the success-competitive arms
sit at $11$--$16$ at $K{=}1$, and only random tie-break approaches it
at $K{=}4$ ($3.2$) at the price of lower success), the lowest
censored time-to-success, and a message cost ($593$--$1050$\,B)
comparable to or below the assignment-style alternatives
(nearest/greedy $423$--$2053$\,B). The occupancy penalty and random
tie-break fail to separate from the unprotected baseline at $K{=}1$.
We report this as a bounded advantage, not dominance.

\begin{figure}[t]
\centering
\includegraphics[width=\linewidth]{figures/fig_warp_drive}
\caption{Warp Drive on unconditional metrics (840 episodes; bootstrap
CIs). \textbf{(a)} Success rate rises at every cadence with disjoint
CIs, and every fast-$K$ protocol arm exceeds the $K{=}8$ sweet-spot
baseline (dashed). \textbf{(b)} Censored time-to-success falls
correspondingly. The claim that the protocol removes the main
argument against fast sharing rests on these unconditional
quantities, not on the (denominator-sensitive) conditional rates.}
\label{p2:fig:drive}
\end{figure}

\begin{table}[h]
\centering
\caption{Unconditional results, 840 episodes (scarcity cells
$\times$ random/asymmetric layouts $\times$ 20 seeds). Success rate
rises under Warp Drive at every cadence with disjoint CIs, censored
time falls, and fast sharing with the protocol beats the $K{=}8$
sweet-spot baseline (Figure~\ref{p2:fig:drive}). Lock census shows the
conditional-metric denominator shrinking --- which is why the claim is
stated on unconditional quantities.}
\label{p2:tab:wd}
\small
\begin{tabular}{lccc}
\toprule
Arm & success rate & $t_{\text{cens}}$ (fail $=120$) & lock census \\
\midrule
$K{=}1$ base & 0.592 [0.574, 0.608] & 55.3 [53.7, 57.2] & 0.986 \\
$K{=}1$ +WD  & \textbf{0.620} [0.609, 0.629] & 53.1 [52.0, 54.4] & 0.985 \\
$K{=}2$ base & 0.572 [0.549, 0.592] & 57.2 [55.1, 59.7] & 0.973 \\
$K{=}2$ +WD  & \textbf{0.614} [0.601, 0.624] & \textbf{53.2} [52.0, 54.6] & 0.926 \\
$K{=}4$ base & 0.577 [0.554, 0.598] & 57.3 [55.1, 59.8] & 0.944 \\
$K{=}4$ +WD  & \textbf{0.620} [0.609, 0.630] & \textbf{52.2} [50.9, 53.7] & 0.739 \\
$K{=}8$ base & 0.583 [0.558, 0.604] & 55.2 [52.9, 57.9] & 0.682 \\
\bottomrule
\end{tabular}
\end{table}

On the conditional axis, recovery is measured against the scarcity
ceiling $\min(m,T)/m$ (no protocol conjures extra water): $0.90 / 2.15
/ 6.6$ of the recoverable collapse at $K{=}1/2/4$; values above 1
arise because reservation suppression removes doomed locks from the
$\mstar$ population. Together with Table~\ref{p2:tab:taxi}, this completes
the arc: raw warp fails at completion (\S\ref{p2:sec:raw}); with a
coordination contract attached, warp at fast cadence becomes the best
regime. The cadence sweet spot of the companion Q/R/M/C study \citep{songlines1} is
reinterpreted: moderate $K$ was never about information freshness ---
it was an implicit rate limiter on warp collisions, and an explicit
contract does the job better.


\section{Cross-session warp on an LLM substrate}
\label{p2:sec:llm}

Finally, the strict warp crosses both an agent boundary and a session
boundary on a language substrate. In session 1, agent A (a local
LLM tag extractor over natural-language observations; Ollama
\method{llama3.1}, deterministic caching, zero API cost) sweeps a
three-room text household and its symbolic memory is dumped to disk.
In session 2, agent B --- a fresh agent in a fresh environment
instance, with the stock directional oracle removed so that long-range
knowledge can only come from memory --- loads A's dump as foreign
provenance and solves ``bring the apple to the table.'' Every lock B
makes on the apple or the table is $\varphi{=}1.0$ strict $\wstar$;
execution is layered (the LLM makes the semantic decisions --- which
remembered place to pursue, when to take/put --- while locomotion to
the locked waypoint is a deterministic motor primitive, matching the
substrate's planner stack).

\begin{table}[h]
\centering
\caption{Cross-session warp over 10 randomised layouts. Completion is
attributed at the environment's affordance radius (take/put act within
Manhattan distance 1, and the witness tags anchors on within-reach
cells).}
\label{p2:tab:llm}
\small
\begin{tabular}{lccc}
\toprule
 & success & mean $t_{\text{succ}}$ & layouts with completed strict $\wstar$ \\
\midrule
warp (A's memory) & \textbf{10/10} & 8.3 & \textbf{10/10} \\
control (empty memory) & 0/10 & --- & --- \\
\bottomrule
\end{tabular}
\end{table}

This experiment provides a controlled demonstration of cross-session,
cross-agent strict semantic warp on an LLM substrate. The $10/10$
versus $0/10$ contrast shows that foreign memory is necessary for completion in this controlled task.

\paragraph{Scaling to a live collective.} With $N{=}3$ LLM agents
sharing by peer broadcast under scarcity ($T{=}2$ apples, four
layouts, $K\in\{2,8\}$), the full warp phenomenology reproduces on the
language substrate (all predictions registered before the runs):
strict $\wstar$ locks occur; warp share falls with cadence
($0.417$ at $K{=}2$ vs.\ $0.083$ at $K{=}8$ --- the same shape as the
grid $0.554\!\to\!0.105$ and VMAS $0.723\!\to\!0.000$); warp
collisions appear in $7/8$ episodes. A paired comparison against the
no-communication baseline on identical layouts shows raw sharing
\emph{lowering} success from $2/3$ to $1/3$ per episode --- the
collision cost exceeds the information value on the LLM substrate
exactly as in \S\ref{p2:sec:gain}. Attaching the unchanged Warp Drive
protocol to the LLM collective provides a corresponding language-agent evaluation: pooled
success rises from $0.333$ to $0.417$ at both cadences with rollback
latencies again ${\approx}K$ ($1$ tick at $K{=}2$, $7$ at $K{=}8$).
The recovery is partial (${\approx}25\%$ of the collapse, vs.\
$90$--$100\%$ on the grid): on a language substrate the collision
pathology is compounded by semantic-level losses (hallucinated anchor
tags, noisy locks) that no reservation can cure --- a measurable
decomposition we leave to future work. The admissibility boundary also survives LLM-produced
evidence: with the gate constant $c$ set to the confidence the
extractor actually emitted ($c_{\text{meas}}{=}0.900$), all four
registered breakpoints land within one tick of prediction.

Two interface findings transfer to any LLM-over-symbolic-memory
design: the LLM query former free-generates retrieval tags outside the
memory's vocabulary (requiring an any-of gate), and small local models
are reliable at semantic commitment but not at coordinate arithmetic
or systematic exploration (requiring the layered execution above and
an explicit exploration tier --- without one, discovery never happens
and the warp has nothing to feed on).

\begin{figure}[t]
\centering
\includegraphics[width=0.62\linewidth]{figures/fig_warp_universality}
\caption{One quantity, three substrates: warp share
$P(\wstar|\mstar)$ vs.\ broadcast cadence $K$ on the symbolic grid,
the continuous VMAS scenario, and the LLM text-world. Across the
three tested substrates, warp share exhibits a qualitatively similar
cadence-dependent pattern; the small LLM arm (8 episodes, no CIs)
provides directional evidence only.}
\label{p2:fig:universality}
\end{figure}

Figure~\ref{p2:fig:universality} overlays the warp-share--cadence curve
across all three substrates: the quantity behaves as an
architecture-level pattern that remains qualitatively consistent
across the three tested substrates.


\section{From target commitments to relational routes}
\label{p2:sec:bridge}

This first movement introduced semantic warp as a provenance-conditioned
annotation of a memory-based commitment. The annotation makes it
possible to distinguish direct observation from peer-supported action,
to estimate the counterfactual contribution of foreign memory, and to
derive an explicit temporal boundary for its authority under a
trust--staleness gate.

The experiments show that foreign evidence should not be evaluated
only by the direct completion rate of foreign locks. Such locks are
fragile under scarce targets, but they can provide substantial
downstream value by transporting recipients into self-confirmation
range. The same analysis localizes the principal cost of fast sharing
to duplicated commitments rather than to evidence quality. A minimal
reservation-and-rollback protocol addresses this failure and improves
unconditional success across the tested fast cadences.

The resulting conclusion is narrower than a general claim that shared
memory is beneficial. Foreign evidence becomes useful when its source,
age, and authority are represented explicitly and when its use is
coupled to a coordination protocol. The remaining sections extend this
result from isolated target commitments to provenance-bearing route
structures and remove the shared-frame assumption.

\paragraph{Outlook: the contract as one function among several.}
The contract perspective suggests a reading we state as a
falsifiable program rather than a conclusion: memory is not one
mechanism but several functions with distinct failure signatures.
The trust$\times$staleness gate studied here implements
\emph{memory-for-trust} (ablate it and completion collapses through
stale and poisoned commits --- the trust-flip and the raw-exchange
negative are exactly that signature); it is distinct from
\emph{memory-for-return} (own consolidation; its ablation starves
candidate generation) and from \emph{memory-for-telling} (what is
worth broadcasting at all; its ablation leaves receivers'
retrieval unimproved on tags they undersample). Under this reading
the right validity criterion for shared memory is not direct
verification of content but \emph{ecological viability}: whether a
commitment formed on foreign evidence can still be carried to
completion under the contract --- which is precisely the quantity
$P(\cstar \mid \wstar)$ and its closed-form boundary make measurable.

\FloatBarrier



\section{Routes and meaning-based identity: the transferable unit}
\label{p3:sec:routes-lead}
\emph{What transfers between agents is the song: edge-provenanced routes, landmark-anchored place identity without shared frames, the vocabulary ablation, and the song anatomy.}

Transferring the location of a target does not necessarily transfer an
executable navigation strategy. In environments with obstacles,
hazards, or private coordinate frames, a recipient may know where a
target is believed to be while lacking both the route required to
reach it and a reliable correspondence between the sender's places and
its own observations. We therefore model the transferable memory unit
as a provenance-bearing relational route. Semantic landmarks identify
corresponding places, route edges encode transport between them, and
edge-level provenance records which parts of the route are known only
through another agent. A strict route-warp event is defined as the
execution of a connected foreign subpath that the recipient has not
previously traversed. Experiments show a clear transition between
place-only and route-level transfer: isolated target records fail at
the first obstructing structure and can perform worse than unguided
exploration, whereas transferred routes remain executable. Because
route edges are observed at different times, their trust-weighted
evidence decays non-uniformly. The resulting rupture boundary predicts the
first inadmissible edge in all $110$ tested regimes and with zero
index error on $21$ intermediate ruptures. Semantic landmark
constellations recover place correspondence across translations and
quarter-turn rotations, while ambiguous configurations are rejected
rather than materialized as phantom targets. A representation
ablation further separates the two informational functions of the
route: landmarks provide identity and edges provide transport; only
their composition reaches an unseen target. In the tested codec, this
composition uses $6.4\%$ of the bits of a full snapshot before protocol
metadata. The results support relational, provenance-conditioned route
transfer as the appropriate unit for executable foreign memory. The
study remains limited to controlled semantic observations and discrete
frame transformations; continuous pose estimation and open-ended
perceptual landmark discovery are outside its scope.

\FloatBarrier
\subsection{From target records to executable routes}
\label{p3:sec:intro}

Collective navigation memory is commonly represented as a set of
places, coordinates, or target descriptors. Such representations can
support target retrieval, but they omit two requirements of executable
transfer. First, the recipient needs a transport structure that
specifies how the relevant places are connected. Second, the sender
and recipient must determine whether their local records refer to the
same physical places, even when their coordinate frames differ. A
shared target label or a nearby coordinate is insufficient when either
of these conditions is absent.

This section addresses both requirements by defining a song as a
relational route with semantic landmarks, ordered transitions, and
edge-level provenance. The term is used as a technical representation:
landmarks provide candidate correspondences between local memories,
whereas route edges specify the sequence of transitions that makes a
target operationally reachable. Provenance indicates which route
segments were acquired by the sender and remain unverified by the
recipient.

The first question is whether route structure provides measurable
value beyond place-only transfer. We evaluate this through controlled
mazes in which the target coordinate is correct but the direct path is
blocked. The comparison is conditional on the recipient's capability:
the traveler here has only local observation and no map or global
replanner, so it cannot reconstruct a feasible path from a destination
record alone. A recipient equipped with a strong planner over a known
map could make a bare coordinate sufficient; the place-only cliff is a
statement about what must be \emph{transferred} when path
reconstruction is not locally available, and intermediate recipients
(local replanning, partial maps, exploration budgets) form a spectrum
we do not sweep here. The second question concerns temporal validity. Because the
observations supporting different route edges have different ages,
the transferred path can remain admissible only up to a finite prefix.
We derive a rupture index by replaying the trust--staleness gate along
the route. The third question is referential: can agents recover the
same route when their private frames differ? We use semantic
constellations around landmarks to generate mutually consistent
correspondences and reject ambiguous matches.

The route-level analysis (contributions 4--7 above) proceeds as a
roadmap: route-level transfer and the place-only cliff
(\S\ref{p3:sec:route}), the edge-specific rupture boundary
(\S\ref{p3:sec:rupture}) and hazard preservation
(\S\ref{p3:sec:hazard}), meaning-based correspondence across private
frames (\S\ref{p3:sec:identity}), and the identity/transport
decomposition with its storage cost (\S\ref{p3:sec:anatomy}).

The central claim is therefore not that relational or topological
navigation is new in itself. The contribution is the composition of
route structure, provenance, temporal validity, and semantic place
correspondence into a measurable transfer object. The present
experiments use controlled semantic vocabularies and discrete
translations and quarter-turn rotations. Arbitrary continuous
$SE(2)$ estimation, rich visual observations, and learned landmark
proposal remain open.


\section{Route-transfer setting}
\label{p3:sec:setting}

Grid worlds ($12{\times}10$ to $14{\times}12$) with semantic water
targets, hazards, and walls; agents with symbolic place memory,
peer broadcast on cadence $K$, and the common memory-grounded planner
used in all experiments
(target lock $=$ the Q/R/M/C $\mstar$ event; completion $\cstar$).
The trust$\times$staleness inclusion rule defined in the companion collective-memory study \citep{songlines3} is used throughout with unchanged constants: evidence
mass $\text{trust}\cdot e^{-\alpha\,\text{age}}\cdot m$, threshold
$\tau{=}0.30$, $\alpha{=}0.05$. Warp events follow the target-level
provenance analysis of Sections~\ref{p2:sec:setting}--\ref{p2:sec:llm}: provenance $\varphi$ frozen at lock time; completion attributed
to a continuous lock from a foreign commit ($\varphi{=}1$ at first
lock) to success.


\section{Route-level transfer}
\label{p3:sec:route}

\paragraph{Edge provenance.} Broadcast snapshots additionally carry the
sender's OWN traversed edges $(a,b,\text{count},\text{last tick})$;
foreign edges are never re-broadcast (the same privacy contract as
place snapshots: no agent reads another's graph, transport only). For
the path $P$ a planner commits,
\[
\phiroute(P) \;=\;
\frac{\sum_{e\in P,\;\pi(e)\neq\mathrm{self}} w(e)}
     {\sum_{e\in P} w(e)},
\qquad w(e) = \text{trust}\cdot e^{-\alpha\,\text{age}(e)}\cdot
\log(1{+}\text{count}(e)),
\]
the same weights the route gate uses. A \textbf{strict $\rwstar$} is a
committed, traversed, connected sub-path of $\ge 5$ edges with
$\phiroute = 1$: the whole verse is someone else's. Place warp is the
degenerate case (path length $0$: only the target is foreign), giving
a two-dimensional provenance plane $(\varphi_{\mathrm{place}},
\phiroute)$.

\paragraph{First strict $\rwstar$.} On a comb maze with measured
detour $D{=}2.65$, a traveler commits and walks a fully foreign
$45$-edge route to completion in $54$ ticks; with the same foreign
PLACE evidence but no edges, the standard planner never arrives
(300-tick budget), and blind exploration takes $126$.

\subsection{Performance transition under route transfer}
\label{p3:sec:cliff}

\begin{figure}[t]
\centering
\includegraphics[width=\linewidth]{figures/fig_route_cliff}
\caption{The cliff (paired arms on identical mazes, 20 seeds per
measured-detour bucket). \textbf{(a)} Place-only transfer does not
degrade with detour --- it collapses at the FIRST wall, located by the
probe bucket between $D{=}1.00$ and $D{=}1.38$; route transfer
completes everywhere. \textbf{(b)} Below the cliff, knowing WHERE
without HOW is \emph{worse than knowing nothing}: the known-target
manhattan pull pins the planner against walls ($-150$ to $-168$ ticks
vs.\ blind exploration), while in the open world the same knowledge is
worth $+57$.}
\label{p3:fig:cliff}
\end{figure}

\begin{table}[h]
\centering
\caption{Three arms on identical mazes (paired by seed), bucketed by
MEASURED detour factor $D$; 20 seeds per bucket, censored gains
(fail $=300$). Place-only transfer does not degrade with $D$ --- it
collapses at the first wall, and is worse than blind exploration
there (the known-target manhattan pull pins the planner against
walls). Route transfer completes everywhere; every completion is a
strict $\rwstar$ with $\phiroute=1$.}
\label{p3:tab:cliff}
\small
\begin{tabular}{lcccc}
\toprule
$D$ & gain route-vs-place [CI] & gain route-vs-blind [CI] & succ r/p/b \\
\midrule
1.00 & $+4.7$ $[+3.4,+6.0]$ & $+62$ $[+44,+80]$ & 20/\textbf{20}/20 \\
1.38 & $+281$ $[+279,+283]$ & $+113$ $[+77,+155]$ & 20/\textbf{0}/16 \\
2.03 & $+273$ $[+271,+275]$ & $+123$ $[+91,+160]$ & 20/\textbf{0}/17 \\
3.32 & $+262$ $[+259,+264]$ & $+110$ $[+72,+152]$ & 20/\textbf{0}/15 \\
\bottomrule
\end{tabular}
\end{table}

Figure~\ref{p3:fig:cliff} shows both faces of the result.
Two pre-registered predictions failed here, in instructive ways.
The negative control at $D{=}1$ was registered in ticks
($|\text{gain}|\le 3$) and came out $+4.7$: the gap is \emph{turn
economy} (BFS routes walk straight legs; the dominant-axis planner
alternates axes), not route knowledge --- in path LENGTH the gap is
exactly $0.0$ $[0.0,0.0]$ across all 20 open-grid cells. And ``gain
grows with $D$'' is unmeasurable: with place success at $0$ the gain
saturates at the censoring ceiling. The correct form of the relationship is a
step located by the probe bucket at $D\in(1.00,1.38)$, i.e.\ the
first wall.


\section{The rupture boundary: edge-wise route validity}
\label{p3:sec:rupture}

The witness stamps edge $i$ at tick $i$: its trace ages
non-uniformly, the start of the route always older than the end, and
a traveler delayed by $a_0$ \emph{chases a fading song}. Under the
standard gate the binding edge is the CURRENT one (the oldest
remaining), giving three regimes on the $a_0$ axis: complete;
rupture at a predictable path index; dead on arrival. The predictor
is fully deterministic --- a no-decay calibration run records the
traveler's kinematics (path index per tick, turns included), and the
registered prediction replays the closed-form gate over that
trajectory, the same deterministic replay procedure used for the
$6/6$ hold-out in Section~\ref{p2:sec:law}.

Across 5 mazes $\times$ 11 delays $\times$ 2 trusts (110 cells,
registered before execution): \textbf{110/110 regime matches; all 21
mid-route ruptures at exactly the predicted edge index (error $0$)};
and a route trust-flip --- at trust $0.7$ the single-traversal edge
mass gives $\mathrm{age}_{\max} = 9.6 < L$, so \emph{no route ever
commits} (0 commits observed; the discrete off-switch of the
warp admissibility boundary of Section~\ref{p2:sec:law}, now for paths). The rupture of a transferred route is
thus not noise but a contract-level regularity: the same constants
$(\alpha,\tau,\text{trust})$ that set WHEN a place-warp gate closes
set WHERE a route-warp breaks.


\section{Hazard preservation along the valid route prefix}
\label{p3:sec:hazard}

\begin{table}[h]
\centering
\caption{Open fields dense with passable-but-punished hazards
(density $0.20$), the witness's route hazard-avoiding by experience;
20 paired worlds, decay off (the safety axis isolated from the aging
axis). Both transfer modes COMPLETE --- the wall axis
(Table~\ref{p3:tab:cliff}) and the hazard axis are independent; they
differ in the price.}
\label{p3:tab:hazard}
\small
\begin{tabular}{lccc}
\toprule
Arm & hazard hits [CI] & completion & mean $t$ \\
\midrule
route (foreign safe path) & \textbf{0} (all 20 cells) & 20/20 & 16.6 \\
place (WHERE only) & 3.25 $[2.45, 4.10]$ & 20/20 & 15.9 \\
blind & 28.7 & --- & --- \\
\bottomrule
\end{tabular}
\end{table}

\begin{figure}[t]
\centering
\includegraphics[width=0.92\linewidth]{figures/fig_route_hazard}
\caption{The song protects. \textbf{(a)} Hazard hits per episode: the
traveler on the witness's route preserves the witness's measured risk profile within the validated route prefix ($0$ in
all 20 worlds); the place-only traveler pays $3.25$ hits for the same
completion; blind exploration pays $28.7$. \textbf{(b)} Composition
with the rupture boundary: in the rupture regime the traveler takes zero
hits while the song lasts, and the risk resumes after the break ---
which occurs at the predicted rupture index in all $19/19$ cells.}
\label{p3:fig:hazard}
\end{figure}

The traveler following the witness's path preserves the witness's
measured risk profile within the validated route prefix, at a time premium of $+0.7$ ticks
(Figure~\ref{p3:fig:hazard}a). Composing with \S\ref{p3:sec:rupture}: in the rupture regime (19/20 worlds) the
traveler takes \textbf{0 hits before the rupture} and $0.58$
$[0.11,1.32]$ after it, when it falls back on place knowledge and
re-enters the field --- with all 19 ruptures again at the
predicted rupture index (error $0$; the edge-exact predictor replicated on
a new world class with no adjustment). The song protects while it
lasts, and the measured protection ends at the edge predicted by the rupture boundary.


\section{Meaning-based place correspondence}
\label{p3:sec:identity}

Everything so far in this paper rides
on a shared coordinate frame. We now remove it: each agent's private
frame is the true grid under a secret translation (later: plus a
secret $90^\circ$-multiple rotation), and broadcasts carry private
coordinates that are meaningless to the receiver until it decides
what the sender's places MEAN.

\paragraph{Two levels of `the same', and what a landmark is.}
Cross-agent place identity has two levels that must not be conflated.
The \emph{semantic predicate} says two cells belong to the same
desired class (\texttt{water\_source}, GOOD): it travels in every
snapshot, and it is what task satisfaction is defined over. But no
predicate can decide \emph{entity correspondence} --- that two memory
entries denote the same world referent --- because a predicate is
exactly what aliases: all water looks like water. Route-relative
description (`$n$ steps past the last GOOD') does not decide it
either: it is provenance of a path, not identity of its endpoints
(\S\ref{p3:sec:route}--\S\ref{p3:sec:rupture} measure what it buys and
where it breaks). What decides correspondence in this series is a
\emph{landmark}: a place whose local semantic constellation supports
a mutually-unique match across agents. Landmark is a \emph{role}, not
a cell type --- which tag classes can play it is itself measured
below (\S\ref{p3:sec:ablation}).

\paragraph{Fingerprints and the handshake.} A place fingerprint is the
local semantic constellation around a visited cell --- which content
tags (hazards, walls, water) sit at which relative offsets ---
translation-invariant and coordinate-free. Matching is
\emph{mutually unique} (a pair is accepted only if each side is the
other's sole candidate above an exact-equality threshold; the sender
knows far more places than the receiver, and one-directional
uniqueness aliases distant look-alikes into the receiver's
neighborhood), fingerprints need $\ge 2$ salient keys, and the frame
translation is the modal delta of $\ge 3$ agreeing pairs. Landmarks
must carry at least one CONTENT tag: border-only (`void') patterns
describe the shape of the world, not a place, and are excluded ---
the lesson of the corner trap below.

\paragraph{Witness--traveler results.} Over 4 offsets $\times$ 2
hazard densities $\times$ 10 seeds: the semantic aligner recovers the
secret offset \emph{exactly in 100\%} of landmark-world runs
(mean $15.1$ mutual pairs) and completes with strict $\wstar$
($\varphi{=}1$) every time; the coordinate baseline under any nonzero
offset locks a non-water cell in $100\%$ of runs (its rare
`successes', $6.7\%$, are the displaced route accidentally crossing
the true water). In the aliased world (featureless interior band) the
two identity regimes break in OPPOSITE directions: semantic identity
makes \emph{zero} locks --- it fails closed --- while coordinate
identity keeps confidently locking wrong cells --- it fails open.
This asymmetry, not raw accuracy, is the deepest argument for
meaning-based identity: when the world does not support
identification, the semantic system knows it.

\subsection{Integrated correspondence and route-transfer stack}
\label{p3:sec:stack}

\paragraph{Online alignment in the peer stack.} With $N{=}4$ agents
under scarcity, all in distinct secret frames, alignment develops
DURING the episode: agents explore, co-visited landmarks accumulate,
and per-sender frames lock in mid-run (onset in $100\%$ of episodes,
mean tick ${\approx}15$). Pooled success: shared-frame oracle
$0.456$, coordinate under misalignment $0.350$, semantic $0.444$ ---
\textbf{88\% of the gap recovered}, with strict $\wstar$ functioning
through recovered frames.

\paragraph{Validity boundary under recovered frames.} Re-running the warp
admissibility boundary of Section~\ref{p2:sec:law} with the trust$\times$staleness gate ON TOP of semantic
identity under a frame offset gives breakpoints $20/14/9$ ---
\emph{the same integers as the shared-frame closed form, exactly}.
The gate and the frame recovery are orthogonal layers: the boundary does
not notice that the shared grid is gone.

\paragraph{No shared north: translations and quarter-turn rotations.} With secret rotations added,
one rotation hypothesis per $90^\circ$ reuses the translation aligner
and the winner must \emph{dominate} ($\ge 2\times$ the runner-up's
consensus). Results: exact $(r^\ast,\delta^\ast)$ recovery in
\textbf{80/80} rotation$\times$offset$\times$seed combinations with
$100\%$ strict-$\wstar$ completion; in a $180^\circ$-symmetric world
--- where the frame is unrecoverable \emph{in principle} --- the
aligner refuses ($0/10$ locks) while coordinates fail open ($10/10$
wrong locks); and the admissibility boundary again does not notice
($\mathrm{bp}=20$ under a $90^\circ$-rotated witness). The
translation-only aligner, for contrast, never recovers a rotated
frame ($0/60$) but does not always fail closed: in $2/60$ cells a
$180^\circ$-symmetric hazard substructure lets it fail OPEN --- the
empirical motivation for the dominance rule.

\begin{figure}[t]
\centering
\includegraphics[width=0.92\linewidth]{figures/fig_route_identity}
\caption{Meaning-based identity under discrete rigid-frame transformations ($\mathbb{Z}^2\rtimes C_4$: translations and quarter-turn rotations). \textbf{(a)}
Every strict-$\wstar$ lock the semantic stack makes points at real
water ($8/8$); the coordinate stack makes $992$ locks, all phantoms
($0/992$); frame recovery itself is exact in $80/80$
witness--traveler combinations. \textbf{(b)} The failure asymmetry in
worlds where identification is impossible (featureless band; a
$180^\circ$-symmetric layout where the frame is unrecoverable in
principle): semantic identity refuses to lock --- fails closed ---
while coordinate identity keeps confidently locking wrong cells ---
fails open.}
\label{p3:fig:identity}
\end{figure}

\paragraph{The corner trap.} In the full stack, raw success is
non-discriminating under rotations (rotated coordinate poison mostly
maps off-grid and degrades into pseudo-exploration), so the
registered discriminator is \emph{warp-lock precision}: the share of
strict $\wstar$ locks pointing at real water. Before the void filter,
early border fingerprints (agents start in corners; rectangle borders
are $180^\circ$-symmetric) fed wrong-frame hypotheses and semantic
precision was $0.159$; with borders excluded as non-landmarks it is
\textbf{$1.000$ ($8/8$) vs.\ $0.000$ ($0/992$ phantom locks) for
coordinates}, with alignment onset in $100\%$ of episodes and clean
regressions on all earlier experiments
(Figure~\ref{p3:fig:identity}).

\paragraph{A stress surface instead of a unit test.} Because
``$80/80$ exact'' certifies only the constructed worlds, the matcher
was subjected to a $5{,}280$-trial stress sweep ($20$ seeds
$\times$ $4$ transforms per level; artifact
\texttt{tmp/cluster/warp/frame\_stress\_full/}), separating
\emph{fail-open} (wrong transform accepted) from \emph{false-reject}
(coverage cost) throughout. Under per-observation tag noise the
matcher is strictly fail-closed across the entire grid
(FN $\le 0.6$, FP $\le 0.3$, and their products: fail-open $= 0$ in
all $42$ noise levels; degradation is pure refusal, with FP more
damaging than FN --- $0.975$ refusal at FP$=0.3$ vs.\ $0.45$ at
FN$=0.6$). Deleting landmarks ($k$ of $21$) is a coverage cost up to
$k{=}7$ (exact $\ge 0.86$, fail-open $0$), with \emph{rare} fail-open
appearing under heavy deletion ($1/80$ at
$k \in \{8, 9, 13, 14\}$). Deliberately ambiguous worlds (repeated
interior constellations, $C \in \{2,3,4\}$) are refused $240/240$
with the ambiguity signal raised in $100\%$ of trials; a
border-anchored variant leaks $2/80$ (the void-context of the border
acts as an implicit landmark and can betray). The one systematic
fail-open channel is \emph{silent vocabulary mismatch}: under a
permuted tag alphabet the matcher accepts a wrong transform in
$2.5$--$5\%$ of trials ($5$ derangements $\times$ $80$), because
mutual-uniqueness consensus can lock onto a coincidental
correspondence. Synonym renaming with a partial translator, missing
concepts, and one-to-many coarsening all remain fail-closed at every
tested level (monotone refusal, fail-open $0$). The identity claims
are therefore conditional on a shared (or explicitly aligned)
vocabulary --- \S\ref{p2:sec:limits} --- and the symbol-alignment
layer is the natural companion mechanism.

\subsection{Causal contribution of the landmark vocabulary}
\label{p3:sec:ablation}

The identity machinery rides on a fixed landmark vocabulary $\Sigma =
\{\texttt{wall}, \texttt{hazard}, \texttt{water}, \texttt{void}\}$:
the tag classes allowed into a fingerprint. Every claim above is
therefore conditional on a hand-chosen alphabet. The last experiment
makes the alphabet the treatment: witness--traveler frame recovery
(worlds carrying both scattered walls and scattered hazards; $3$
offsets $\times$ $12$ seeds per arm) under seven vocabularies, from
full $\Sigma$ down to single classes.

\begin{table}[h]
\centering
\caption{Landmark-vocabulary ablation ($36$ cells per arm, $252$
total; registered predictions A1--A4, all PASS). Impoverishing the
alphabet costs \emph{coverage} --- recovery converts into refusal ---
but never \emph{correctness}: every transported lock, under every
vocabulary, goes through the exactly-recovered frame. Attribution =
share of consensus-pair key mass under full $\Sigma$.}
\label{p3:tab:ablation}
\small
\begin{tabular}{lcccc}
\toprule
Vocabulary & exact recovery & mean pairs & lock precision & fail-open \\
\midrule
full $\Sigma$ & $1.00$ & $33.2$ & $1.00$ & $0$ \\
$-$ void & $1.00$ & $32.0$ & $1.00$ & $0$ \\
$-$ hazards & $0.92$ & $11.8$ & $1.00$ & $0$ \\
$-$ walls & $1.00$ & $15.5$ & $1.00$ & $0$ \\
hazards only & $1.00$ & $13.7$ & $1.00$ & $0$ \\
walls only & $0.83$ & $10.8$ & $1.00$ & $0$ \\
void only & $0.00$ & $0.0$ & --- & $0$ \\
\bottomrule
\end{tabular}
\end{table}

Four registered predictions, four passes (Table~\ref{p3:tab:ablation}).
\textbf{(A1) The ablation is fail-safe:} in $252$ cells there is not
one fail-open transport, consistent with the fail-closed asymmetry of
\S\ref{p3:sec:identity} holds not just against ambiguous \emph{worlds}
but against impoverished \emph{alphabets}. \textbf{(A2) Void is
scaffolding, not identity:} border patterns carry $0.000$ of the
consensus key mass, and dropping them changes nothing --- consistent
with the corner trap, where they actively poisoned rotation
hypotheses. \textbf{(A3) One content class suffices:} hazards alone
recover every frame; walls alone recover $0.83$ and refuse the rest.
Correspondence needs only $\ge 3$ mutually-unique constellations over
\emph{some} content class --- a bounded-rationality budget, not a rich
ontology. \textbf{(A4) Attribution is measurable:} hazard
constellations carry $0.56$ of the consensus mass, walls $0.44$.
Which features function as landmarks is thus not a modelling
assumption that happened to work --- it is a measurable property of
the world$\times$vocabulary pair, and the machinery degrades along
the safe axis whenever the vocabulary is wrong.

\subsection{Functional decomposition of the route representation}
\label{p3:sec:anatomy}

\S\ref{p3:sec:identity} opened with a claim we have so far only argued:
route-relative description (`$n$ steps past the last GOOD') is
provenance of a path, not identity of its endpoints. The last
experiment measures it. A witness compresses its route to the water
into a purely \emph{relational} song --- a sequence of couplets
$(\sigma, e)$, where $\sigma$ is the landmark constellation signature
and $e$ the exact beat vector from the previous couplet --- carrying
\emph{no global coordinates whatsoever}. A band-limited traveler who
has never seen the water receives it. Four registered arms (bit codec
fixed before the runs; $12$ paired worlds):

\begin{itemize}[leftmargin=16pt,itemsep=2pt,topsep=2pt]
\item \textbf{Beats only} (the pure `$n$ steps from the last GOOD'):
  usable only by \emph{assuming} the anchor. With a shared anchor it
  is exact ($11/11$, the negative control); with anchor disagreement
  it locks every time and is right \emph{never} ($11/11$ fail-open
  mislocks, displaced by exactly the anchor offset). Beats do not
  carry identity --- they parasitise on the identity of the anchor.
\item \textbf{Signatures only}: re-identifies co-visited places but
  has no edges to walk --- \emph{zero} transports of the never-seen
  target. Nodes alone cannot reach the unseen.
\item \textbf{Signatures $+$ beats}: couplets matched mutually-uniquely,
  the beat chain checked for metric consistency between matched
  couplets (a loop-closure test), the target dead-reckoned from the
  last matched landmark. Coverage $1.0$, every transport exact, zero
  fail-open --- at $366$ bits per song.
\item \textbf{Full snapshots} (the machinery of
  \S\ref{p3:sec:identity}): the information upper bound --- same
  correctness at $5{,}690$ bits.
\end{itemize}

All four registered predictions passed. The song is not a lighter
snapshot; it is a different \emph{kind} of object: identity lives in
the nodes (landmarks), transport lives in the edges (beats), and
their composition delivers full-snapshot transport at $6.4\%$ of the
bit cost while inheriting the fail-safe property of
\S\ref{p3:sec:ablation} --- everything that goes wrong goes wrong as a
refusal. (This $6.4\%$ is the \emph{pure song codec}: nodes and
edges only. A deployed collective-memory protocol adds provenance,
timestamps, world versions, and a utility certificate; the
full-protocol payload ratio, measured in the companion collective-memory study \citep{songlines3}, is correspondingly larger.)


\section{Related work}
\label{p2:sec:related}

\paragraph{Shared memory in agent collectives.} Generative
Agents~\cite{park2023generative} and multi-agent LLM frameworks share
memory streams or context but expose no lock interface and no
per-source decomposition of a merged belief, so provenance-conditioned
completion is not measurable there. Vector memories (e.g.\
MemGPT~\cite{packer2023memgpt}) merge by embedding arithmetic;
the merge is not source-decomposable, which blocks both $\varphi$ and
the masked-replay counterfactual.
\paragraph{Emergent communication.} CommNet~\cite{sukhbaatar2016commnet}
and TarMAC~\cite{das2019tarmac} learn continuous channels; the
companion Q/R/M/C study \citep{songlines1} shows such channels collapse stage observability. Warp
requires an explicit materialization event, which these architectures
do not expose.
\paragraph{Stigmergy.} Pheromone-style coordination is an
\emph{anonymous} warp: evidence transfers without provenance. Our
results quantify what anonymity costs (shared-bus warp share is
maximal while its warp completion is minimal) and what
provenance buys (trust gates, reservations, rollback).
\paragraph{Optimistic concurrency.} The rollback primitive is Time
Warp~\cite{jefferson1985} transplanted from parallel simulation to
memory-driven navigation: locks are optimistic computations,
occupancy evidence is a straggler message, anti-$\mstar$ is the
anti-message. We adapt Time-Warp-style optimistic rollback to
decentralized coordination over memory-grounded target commitments.


\label{p3:sec:related}

\paragraph{Map merging and place recognition.} Multi-robot SLAM
treats inter-map data association as an estimation problem over
metric features~\cite{cadena2016slam}; visual place
recognition~\cite{lowry2016vpr} matches appearance. Our contribution
is not a competing matcher but the measurement layer on top:
provenance-conditioned completion through recovered frames, the
fail-closed/fail-open asymmetry as a registered prediction, and
boundary invariance under frame removal.
\paragraph{Route sharing.} Stigmergy transfers paths anonymously
(pheromone gradients); learned communication (e.g.\
CommNet~\cite{sukhbaatar2016commnet}) transfers policies opaquely.
Neither exposes an edge-provenance decomposition, so neither can
state --- let alone predict --- where a transferred route breaks.
The rupture predictor is Time Warp~\cite{jefferson1985} thinking
applied to memory decay: deterministic replay of a gate over
calibrated kinematics.
\paragraph{Songlines.} The narrative framing corresponds to a functional distinction in
both directions: the song is a sequence (route warp), its places are
recognised by meaning (semantic identity), it fades verse by verse
(rupture boundary), and it protects the walker while it lasts (hazard
stratification).


\section{Limitations}
\label{p2:sec:limits}

\paragraph{Evidence classes.} The paper's results occupy different
rungs of a causal-strength ladder, and we mark them explicitly rather
than let the strongest reading spread to all: the $\varphi$-stratified
completion gap is a \emph{descriptive stratified association} (with
the selection caveat of \S\ref{p2:sec:raw}); the lock-chain $51\%$ is
a \emph{descriptive precedence} statistic; the masked-replay warp gain
is a \emph{paired system-level counterfactual} (interference
included); the admissibility and rupture boundaries are \emph{formal
consequences of the implemented contract}, validated by registered
hold-outs against the implementation; the Warp Drive gains are a
\emph{module-level intervention}; and hazard preservation and the
place-only cliff are \emph{controlled empirical comparisons}
conditional on the stated recipient capability.

The provenance analysis is subject to six main limitations. First,
the primary experiments use a shared coordinate frame and therefore
do not solve cross-agent entity correspondence; the route-level
analysis in the following sections addresses this assumption with
semantic landmark matching. Second, the evidence
is based mainly on controlled grid environments, with a directional
VMAS study and a small text-world LLM demonstration rather than a
photorealistic benchmark. Third, the sign of counterfactual warp gain
is regime-dependent and should not be summarized as a universal
benefit of foreign memory. Fourth, the closed-form admissibility boundary is
specific to the explicit exponential trust--staleness gate and its
pruning schedule. Fifth, the co-lock-pressure analysis is exploratory
and is reported as mechanism evidence rather than as a confirmed
predictive law. Sixth, the LLM experiment uses a limited model and task
family. All deviations from the registered design and all failed
predictions are retained in Appendix~\ref{p2:app:deviations}.


\label{p3:sec:limits}

The route-transfer experiments use controlled navigation substrates
and a predefined semantic vocabulary. Place correspondence is tested
under translations and quarter-turn rotations; this setting should be
described as discrete frame recovery rather than arbitrary continuous
$SE(2)$ estimation. The semantic fingerprints are sufficiently rich
to support exact correspondence in the reported environments and
degrade fail-closed under tag noise, landmark deletion, and
constructed ambiguity (\S\ref{p3:sec:identity}); the measured
exception is silent vocabulary mismatch, where a permuted tag
alphabet induces $2.5$--$5\%$ wrong-transform acceptance --- so the
fail-closed guarantee is conditional on a shared or explicitly
aligned ontology, and open-world appearance variation, perceptual
uncertainty, and learned symbol alignment remain unresolved. The rupture boundary is exact for the
specified edge-wise trust--staleness mechanism and does not by itself
provide a universal safety guarantee. The claim concerning risk is
therefore restricted to preservation of the witness's measured hazard
profile within the validated route prefix. Finally, the $6.4\%$ figure
refers to the pure route representation; a complete deployment cost
must additionally include provenance, timestamps, correspondence
hypotheses, and coordination metadata.


\section{Conclusion}

Taken together, the two movements of this paper give a single answer to
what makes foreign navigation memory executable. The first showed that a
peer-supported \emph{commitment} carries value only under a provenance,
validity, and coordination contract --- direct completion understates
it, a trust--staleness horizon bounds it, and a reservation protocol
makes fast sharing net-positive. The second identifies a
provenance-bearing relational route as the
transferable unit of executable navigation memory. A target record by
itself can identify a destination without specifying a feasible path,
and an edge sequence without a reliable landmark correspondence can
transport the agent toward the wrong referent. The experiments show
that these functions are distinct: semantic landmarks provide
identity, route edges provide transport, and only their composition
supports transfer to an unseen target.

The route representation also exposes a temporal boundary. By
applying the trust--staleness gate to each edge, the method predicts
the route prefix that remains admissible and the edge at which the
transfer should be rejected or replanned. In the controlled studies,
this rupture index matches the observed failure edge exactly under the
implemented trust--staleness gate and the tested route dynamics.
Meaning-based landmark correspondence further removes the shared
coordinate assumption for discrete translations and quarter-turn
rotations and fails closed when the observations are ambiguous.

The results support a precise interpretation of foreign navigation
memory: it must preserve referential identity, transport structure,
provenance, and a validity boundary together. How such memory is
\emph{formed}, consolidated, and admitted into a collective over long
horizons --- with explicit merge, trust, and staleness rules --- is the
subject of the companion collective-memory study \citep{songlines3}.


\clearpage
\appendix

\section{Deviations from the pre-registered design}
\label{p2:app:deviations}
\begin{enumerate}[leftmargin=20pt,itemsep=2pt,topsep=2pt]
\item Reservation tie-break extended to (lock tick, \emph{distance at
  lock}, agent id): a pure id tie-break hands simultaneous locks to a
  distant agent over the target's near-discoverer.
\item Protocol recovery is measured against the scarcity ceiling
  $\min(m,T)/m$ rather than the $K{=}8$ reference, separating
  structural from coordination collapse; unconditional metrics are
  reported alongside (Table~\ref{p2:tab:wd}).
\item Completion breakpoints use the discrete-time refinement of
  \S\ref{p2:sec:law} (boundary constants unchanged; hold-out registered).
\item Collision pressure is operationalized as co-lock count rather
  than co-recipients (structurally constant $N{-}2$ under broadcast).
\item LLM demo: layered execution (LLM semantics, deterministic
  locomotion); completion attributed at the affordance radius.
\end{enumerate}


\section{Reproducibility}
\label{p2:app:repro}
All experiments run on CPU in minutes-to-tens-of-minutes; the LLM demo
uses a local Ollama model with deterministic caching.
\begin{verbatim}
PYTHONPATH=. python experiments/warp/exp_warp_anchor.py
PYTHONPATH=. python experiments/warp/exp_warp_gain.py --seeds 20
PYTHONPATH=. python experiments/warp/analyze_warp_chains.py
PYTHONPATH=. python experiments/warp/exp_warp_gain_individual.py
PYTHONPATH=. python experiments/warp/exp_warp_age_law.py
PYTHONPATH=. python experiments/warp/exp_warp_age_law_holdout.py
PYTHONPATH=. python experiments/warp/exp_warp_drive.py --seeds 20
PYTHONPATH=. python experiments/warp/exp_warp_cross_session.py --layouts 10
\end{verbatim}


\label{p3:app:repro}
\begin{verbatim}
PYTHONPATH=. python experiments/warp/exp_route_warp_r0.py
PYTHONPATH=. python experiments/warp/exp_route_warp_r1.py --seeds 20
PYTHONPATH=. python experiments/warp/exp_route_warp_r2.py --seeds 5
PYTHONPATH=. python experiments/warp/exp_route_warp_r3.py --seeds 20
PYTHONPATH=. python experiments/warp/exp_warp_semantic_identity.py
PYTHONPATH=. python experiments/warp/exp_warp_semantic_stack.py --seeds 20
PYTHONPATH=. python experiments/warp/exp_warp_rotation.py --seeds 10
PYTHONPATH=. python experiments/warp/exp_warp_landmark_ablation.py
PYTHONPATH=. python experiments/song_grammar/exp_s0_song_smoke.py
\end{verbatim}
\FloatBarrier


\section{Full counterfactual sign map}
\label{p2:app:signmap}
\FloatBarrier

\clearpage

\section{Registered predictions and revisions}
\label{p3:app:registered}
All predictions were written to disk before execution; the artifact
retains both versions wherever an operationalisation was revised.
Revisions: R1 negative control re-measured in path length; R1 gain
monotonicity re-formulated as cliff detection; W9 full-stack
discriminator changed from success to warp-lock precision; W8
strict-$\wstar$ functioning extended to include warp-assisted own
completions (taxi chains). The vocabulary ablation (A1--A4,
\S\ref{p3:sec:ablation}) and the song-anatomy experiment (S0.1--S0.4,
\S\ref{p3:sec:anatomy}, bit codec registered with the predictions)
passed as registered, with no revisions.


\begin{thebibliography}{99}
\bibitem[Anonymous(2026a)]{songlines1}
Anonymous.
\newblock Q/R/M/C: Stage-decomposed evaluation of memory-grounded agents.
\newblock Under review, 2026.
\bibitem[Anonymous(2026c)]{songlines3}
Anonymous.
\newblock Utility-gated collective semantic memory: Formation, admission, and long-horizon coordination.
\newblock Under review, 2026.
\bibitem[Packer et~al.(2023)]{packer2023memgpt}
C.~Packer, V.~Fang, S.~Patil, K.~Lin, S.~Wooders, and J.~Gonzalez.
\newblock MemGPT: Towards LLMs as operating systems.
\newblock \emph{arXiv:2310.08560}, 2023.
\bibitem[Park et~al.(2023)]{park2023generative}
J.~S. Park, J.~C. O'Brien, C.~J. Cai, M.~R. Morris, P.~Liang, and M.~S. Bernstein.
\newblock Generative agents: Interactive simulacra of human behavior.
\newblock \emph{arXiv:2304.03442}, 2023.
\bibitem[Sukhbaatar et~al.(2016)]{sukhbaatar2016commnet}
S.~Sukhbaatar, A.~Szlam, and R.~Fergus.
\newblock Learning multiagent communication with backpropagation.
\newblock In \emph{NeurIPS}, 2016.
\bibitem[Das et~al.(2019)]{das2019tarmac}
A.~Das, T.~Gervet, J.~Romoff, D.~Batra, D.~Parikh, M.~Rabbat, and
J.~Pineau.
\newblock TarMAC: Targeted multi-agent communication.
\newblock In \emph{ICML}, 2019.
\bibitem[Jefferson(1985)]{jefferson1985}
D.~R. Jefferson.
\newblock Virtual time.
\newblock \emph{ACM Transactions on Programming Languages and Systems},
7(3):404--425, 1985.
\bibitem[Cadena et~al.(2016)]{cadena2016slam}
C.~Cadena, L.~Carlone, H.~Carrillo, Y.~Latif, D.~Scaramuzza,
J.~Neira, I.~Reid, and J.~J. Leonard.
\newblock Past, present, and future of simultaneous localization and
mapping: Toward the robust-perception age.
\newblock \emph{IEEE Transactions on Robotics}, 32(6):1309--1332, 2016.
\bibitem[Lowry et~al.(2016)]{lowry2016vpr}
S.~Lowry, N.~S\"underhauf, P.~Newman, J.~J. Leonard, D.~Cox,
P.~Corke, and M.~J. Milford.
\newblock Visual place recognition: A survey.
\newblock \emph{IEEE Transactions on Robotics}, 32(1):1--19, 2016.
\end{thebibliography}
\end{document}
